% ============================================================
%  Resultados y análisis
%  Solo resultados y su lectura; el desarrollo estadístico y de
%  incertidumbre está en el Anexo C.
% ============================================================
\section{Resultados y análisis}
\label{sec:resultados}

Los resultados provienen de dos fuentes
(Sección~\ref{subsubsec:equipo_syglab}): los registros del equipo completo de
referencia provisto por SYGlab, que ejecutó los ensayos de seguridad con la
versión inicial del software (Sección~\ref{subsec:resultados_syglab}), y las pruebas de la versión
final en el prototipo propio, armado con los componentes descritos en la
Sección~\ref{subsec:metodologia_hw} (Figura~\ref{fig:conexion_red}) y operado
con la aplicación y el hardware reales.
En el prototipo se usaron dos cargas: un soldador de
\SI{40}{\watt} declarados (carga resistiva) y un ventilador (carga con motor). La
evidencia ---fotografías, base de datos, exportación y scripts de prueba--- se
lista en el Anexo~\ref{anx:evidencia}.

Al momento de las pruebas no se había realizado el contraste con el instrumento
patrón previsto (Sección~\ref{subsec:metodologia_hw}), por lo que los
resultados verifican el funcionamiento del sistema y la corrección de sus
veredictos, pero \textbf{no constituyen una validación metrológica}.

% ------------------------------------------------------------
\subsection{Ensayo de consumo}
\label{subsec:resultados_consumo}
% ------------------------------------------------------------

La Tabla~\ref{tab:resultados_consumo} resume las mediciones con carga. La tensión
de red medida varió entre \SIlist{226;231}{\volt} a lo largo de los días de
prueba.

\begin{table}[htbp]
\centering
\caption{Mediciones de consumo en el banco de pruebas}
\label{tab:resultados_consumo}
\renewcommand{\arraystretch}{1.25}
\small\setlength{\tabcolsep}{4pt}
\begin{tabular}{@{} l p{3.6cm} c c c c @{}}
\toprule
\textbf{Carga} & \textbf{Prueba} & \textbf{$V$ [\si{\volt}]} &
\textbf{$I$ [\si{\ampere}]} & \textbf{$P$ [\si{\watt}]} & \textbf{FP} \\
\midrule
Ventilador & Lectura directa
  & 228,7--228,9 & 0,236--0,242 & 52,1--53,4 & 0,96--0,97 \\
Soldador   & Lectura directa
  & 230,1--230,2 & 0,163--0,164 & 37,5--37,7 & 1,00 \\
Soldador   & Final de línea, 4 ensayos
  & 227,7--228,0 & 0,161--0,162 & 36,6--36,8 & 0,99--1,00 \\
Soldador   & Perfil de consumo, 8 muestras
  & 226,2--227,3 & 0,160--0,161 & 36,1--36,4 & 1,00 \\
\bottomrule
\end{tabular}
\end{table}

Para el soldador, configurado como aparato de calefacción de \SI{40}{\watt}, la
tolerancia aplicable es de \SI{\pm 10}{\percent} (\SIrange{36}{44}{\watt}). La
Tabla~\ref{tab:resultados_veredicto} resume el resultado del ensayo de final de
línea.

\begin{table}[htbp]
\centering
\caption{Veredicto del ensayo de consumo del soldador}
\label{tab:resultados_veredicto}
\renewcommand{\arraystretch}{1.25}
\begin{tabular}{@{} l r @{}}
\toprule
\textbf{Parámetro} & \textbf{Valor} \\
\midrule
Potencia nominal declarada        & \SI{40}{\watt} \\
Intervalo admisible (calefacción, \SIrange{25}{200}{\watt}) & \SIrange{36}{44}{\watt} \\
Potencia media medida ($n = 4$)   & \SI{36,70}{\watt} \\
Desviación respecto del nominal   & \SI{-8,3}{\percent} \\
Coeficiente de variación          & \SI{0,22}{\percent} \\
Incertidumbre expandida ($k = 2$) & \SI{0,23}{\watt} \\
\midrule
\textbf{Veredicto del sistema}    & \textbf{OK} \\
\bottomrule
\end{tabular}
\end{table}

\paragraph{Análisis.}
\begin{itemize}
    \item La \textbf{repetibilidad} es muy buena: los cuatro ensayos de final de
          línea difieren en \SI{0,2}{\watt} y el coeficiente de variación es del
          \SI{0,22}{\percent}, del mismo orden que la resolución del medidor.
    \item La \textbf{incertidumbre} estimada (\SI{0,23}{\watt}) es pequeña frente
          al ancho de la banda de tolerancia (\SI{\pm 4}{\watt}), por lo que la
          dispersión del sistema no compromete el veredicto. El desarrollo del
          cálculo está en el Anexo~\ref{anx:estadistica}.
    \item El resultado queda \textbf{cerca del límite inferior} (\SI{36}{\watt}).
          La cláusula~10.1 pide medir a la tensión nominal y el banco mide a la
          tensión de red disponible, sin regularla. El soldador está marcado para
          \SI{230}{\volt}; como en una carga resistiva la potencia varía con el
          cuadrado de la tensión, su consumo a la tensión nominal se estima en
          \SI{37,4}{\watt} (\SI{-6,6}{\percent}), con lo que el veredicto se
          mantiene. En un aparato más cercano al límite, en cambio, la tensión de
          ensayo podría cambiar el veredicto; el Anexo~\ref{anx:estadistica}
          desarrolla este efecto.
    \item Las lecturas del perfil (\SI{36,2}{\watt} en promedio) son
          \SI{0,5}{\watt} menores que las de final de línea, coherente con la
          menor tensión de red en ese momento (\SI{227}{\volt} frente a
          \SI{228}{\volt}).
\end{itemize}

% ------------------------------------------------------------
\subsection{Prueba de punta a punta}
\label{subsec:resultados_e2e}
% ------------------------------------------------------------

La prueba de punta a punta recorrió la aplicación completa sobre el
banco: detección del hardware, alta de productos, ensayos de final de línea y
cíclico, parada de emergencia, cierre de la ventana durante un ensayo, historial,
exportación y API. Encontró seis defectos, que se corrigieron y se volvieron a
verificar en el banco (Tabla~\ref{tab:resultados_e2e}).

\begin{table}[htbp]
\centering
\caption{Defectos encontrados en la prueba de punta a punta y verificación de la
         corrección en el banco}
\label{tab:resultados_e2e}
\renewcommand{\arraystretch}{1.25}
\small
\begin{tabularx}{\linewidth}{@{} l X X @{}}
\toprule
\textbf{Área} & \textbf{Antes} & \textbf{Después de la corrección} \\
\midrule
Consumo (soldador)
  & \SIlist{27,2;33,2;31,6;32,6}{\watt}: FALLA
  & \SIlist{36,8;36,7;36,7;36,6}{\watt}: OK \\
Modo cíclico
  & La interfaz se liberaba tras el primer ciclo
  & Bloqueada hasta el tercer ciclo (\SI{22,1}{\second}) \\
Error durante un ensayo
  & Interfaz bloqueada, sin aviso
  & Aviso visible, interfaz libre en \SI{0,5}{\second}, relés abiertos \\
Fuga sin medidor
  & NO EJECUTADO
  & FALLA (lectura fallida = FALLA) \\
Historial
  & Columnas de corriente y potencia vacías
  & Completas, con la resolución del medidor \\
Cierre durante un ensayo
  & El ensayo seguía unos \SI{4}{\second}
  & Se detiene en \SI{0,4}{\second}, relés abiertos \\
\bottomrule
\end{tabularx}
\end{table}

El defecto más relevante fue la \textbf{subestimación del consumo en un
\SI{27}{\percent}}: el promedio incluía lecturas tomadas mientras el contactor
cerraba (\SI{0}{\watt}, luego \SI{26}{\watt} y recién después \SI{37}{\watt}), y
un aparato en buen estado resultaba rechazado. Se corrigió descartando las
lecturas de estabilización (Sección~\ref{subsec:criterios_ensayo}). Además, la
parada de emergencia cortó el ensayo en \SI{0,6}{\second} con los relés abiertos,
y la exportación generó un archivo Excel con las hojas de historial, consumo y
verificación.

% ------------------------------------------------------------
\subsection{Ensayos en el equipo completo de SYGlab}
\label{subsec:resultados_syglab}
% ------------------------------------------------------------

La base de datos de la versión inicial del software, instalada en el equipo de
SYGlab, registra \textbf{86 verificaciones con patrones} realizadas en 18
jornadas entre 2023 y 2025. En cada verificación el equipo ensaya dos patrones
conocidos ---uno que debe aprobar y otro que debe fallar--- en corte de neutro,
puesta a tierra y tensión resistida, y registra si obtuvo el resultado esperado
(OK) o no (ERROR). La Tabla~\ref{tab:verificaciones_fdl} resume los resultados;
la base registra además tres ensayos de consumo de un producto Indelplas IV~20.

\begin{table}[htbp]
\centering
\caption{Verificaciones con patrones del equipo completo de SYGlab (versión inicial
         del software, 2023 a 2025)}
\label{tab:verificaciones_fdl}
\renewcommand{\arraystretch}{1.25}
\small
\begin{tabular}{@{} l c c c c c @{}}
\toprule
& \multicolumn{2}{c}{\textbf{Patrón que debe aprobar}}
& \multicolumn{3}{c}{\textbf{Patrón que debe fallar}} \\
\cmidrule(lr){2-3}\cmidrule(l){4-6}
\textbf{Ensayo} & \textbf{OK} & \textbf{ERROR} & \textbf{OK} & \textbf{ERROR} &
\textbf{No ensayado} \\
\midrule
Corte de neutro    & 86 & 0  & 76 & 0 & 10 \\
Puesta a tierra    & 42 & 44 & 69 & 7 & 10 \\
Tensión resistida  & 84 & 2  & 69 & 7 & 10 \\
\bottomrule
\end{tabular}
\end{table}

\paragraph{Análisis.}
\begin{itemize}
    \item El corte de neutro y la tensión resistida se comportaron de forma
          consistente: el patrón que debe aprobar dio el resultado esperado en
          el \SI{100}{\percent} y el \SI{98}{\percent} de las verificaciones.
    \item La \textbf{puesta a tierra} fue la medición más sensible: el patrón que
          debe aprobar resultó en ERROR en 44 de las 86 verificaciones. Coincide
          con que el sistema tenía configurada una corrección de calibración de
          puesta a tierra, y con lo observado en el diseño de este trabajo: la
          medición de resistencias de pocos miliohmios depende de la corriente de
          ensayo (al menos \SI{10}{\ampere}), de la resistencia de contacto y del
          método de medición (Anexo~\ref{anx:valores_referencia}). Esto motivó
          exigir en la versión final que la corriente alcance el mínimo de la
          norma para dar por válido el ensayo.
    \item Solo 39 de las 86 verificaciones resultaron completamente correctas en
          los seis controles, lo que muestra el valor de verificar el equipo con
          patrones antes de cada jornada, función que la versión final conserva.
\end{itemize}

% ------------------------------------------------------------
\subsection{Ensayos de seguridad con la versión final}
\label{subsec:resultados_seguridad}
% ------------------------------------------------------------

En la versión final, los ensayos de fuga, puesta a tierra y rigidez dieléctrica
están implementados con los límites de la norma y verificados con hardware
simulado. Su ejecución con hardware real corresponde a la versión inicial en el
equipo de SYGlab (Sección~\ref{subsec:resultados_syglab}): el prototipo propio no
incluye su hardware de medición, y no hay registros de la versión final en
aquel equipo (Tabla~\ref{tab:estado_ensayos}). En el prototipo se comprobó que, sin ese
hardware, el sistema responde de forma segura: informa FALLA y deja los relés
abiertos.

\begin{table}[htbp]
\centering
\caption{Estado de cada ensayo al cierre del trabajo}
\label{tab:estado_ensayos}
\renewcommand{\arraystretch}{1.25}
\small
\begin{tabularx}{\linewidth}{@{} l l X @{}}
\toprule
\textbf{Ensayo} & \textbf{Estado} & \textbf{Falta para validarlo en banco} \\
\midrule
Consumo (final de línea, perfil, cíclico) & Validado en banco & Contraste con patrón \\
Fuga activa y pasiva & Software verificado
  & PZEM-004T de \SI{10}{\ampere} y red de medición IEC~60990 \\
Puesta a tierra & Software verificado
  & PZEM3, transformador de \SI{2,2}{\volt} y verificar $\geq$ \SI{10}{\ampere} \\
Rigidez dieléctrica & Software verificado
  & Transformador de alta tensión y medición de \SI{5}{\milli\ampere} \\
Verificación con patrones & Software verificado & Hardware de los tres ensayos anteriores \\
\bottomrule
\end{tabularx}
\end{table}

% ------------------------------------------------------------
\subsection{Calidad del software}
\label{subsec:resultados_software}
% ------------------------------------------------------------

Las pruebas automáticas pasaron de 63 a 158 a lo largo del período y cubren los
veredictos de la norma, las tolerancias de consumo, los ensayos con hardware
simulado, los relés (incluida la reconexión), las correcciones de calibración, las
credenciales y todas las pestañas de la interfaz. La integración continua, que no
ejecutaba las pruebas desde el primer commit por un error de configuración, quedó
en verde en sus tres etapas: pruebas, análisis estático y prueba en Docker.

% ------------------------------------------------------------
\subsection{Tiempo de ensayo e intervención del operario}
\label{subsec:resultados_tiempos}
% ------------------------------------------------------------

Como no se realizó una medición comparativa con un operario ensayando con
instrumentos convencionales, la Tabla~\ref{tab:tiempos} presenta una
\textbf{estimación} por aparato. Para el método manual se considera que cada
ensayo exige conectar el instrumento, configurarlo, ejecutar y anotar el
resultado; para el sistema, la duración configurada de cada ensayo
(\SI{5}{\second}) más la conmutación, coherente con lo observado en el banco
(tres ciclos de consumo en \SI{22,1}{\second}).

\begin{table}[htbp]
\centering
\caption{Estimación del tiempo y de las intervenciones del operario por aparato}
\label{tab:tiempos}
\renewcommand{\arraystretch}{1.25}
\small
\begin{tabularx}{\linewidth}{@{} X r r @{}}
\toprule
\textbf{Etapa} & \textbf{Manual} & \textbf{Sistema} \\
\midrule
Conexión del aparato y configuración          & \SI{75}{\second} & \SI{40}{\second} \\
Consumo                                       & \SI{60}{\second} & \multirow{5}{*}{\SI{30}{\second}} \\
Fuga activa                                   & \SI{45}{\second} & \\
Fuga pasiva                                   & \SI{45}{\second} & \\
Puesta a tierra                               & \SI{45}{\second} & \\
Rigidez dieléctrica                           & \SI{60}{\second} & \\
Registro de resultados                        & \SI{60}{\second} & automático \\
\midrule
\textbf{Tiempo total por aparato}             & \textbf{\SI{6,5}{\minute}}
                                              & \textbf{\SI{70}{\second}} \\
\textbf{Intervenciones del operario}          & \textbf{21} & \textbf{4} \\
\bottomrule
\end{tabularx}
\end{table}

Con estas estimaciones, el sistema reduce el tiempo de ensayo en torno al
\SI{80}{\percent} y las intervenciones del operario de 21 (conectar, configurar,
iniciar y anotar en cada uno de los cinco ensayos, y desconectar) a 4 (conectar,
elegir el producto, iniciar y desconectar), es decir, un \SI{81}{\percent}.
Ambos valores deben confirmarse con mediciones cuando estén operativos los cinco
ensayos.

% ------------------------------------------------------------
\subsection{Evaluación de los criterios de éxito}
\label{subsec:resultados_criterios}
% ------------------------------------------------------------

\begin{table}[htbp]
\centering
\caption{Evaluación de los criterios de éxito}
\label{tab:criterios_exito}
\renewcommand{\arraystretch}{1.25}
\small
\begin{tabularx}{\linewidth}{@{} l l X @{}}
\toprule
\textbf{Criterio} & \textbf{Meta} & \textbf{Resultado} \\
\midrule
Tiempo de prueba & $\geq$ \SI{30}{\percent} menos
  & Cumplido según la estimación: alrededor de \SI{80}{\percent} menos
    (Tabla~\ref{tab:tiempos}); a confirmar con mediciones \\
Intervención del operario & $\geq$ \SI{80}{\percent} menos
  & Cumplido según la estimación: de 21 a 4 intervenciones
    (\SI{81}{\percent} menos); a confirmar con mediciones \\
Seguridad & Alarmas y falla segura
  & Cumplido en el banco: relés abiertos ante error, parada y cierre;
    lectura fallida = FALLA \\
Escalabilidad & Adaptable a productos y normas
  & Cumplido: límites por atributos del producto y umbrales configurables \\
Cumplimiento normativo & Límites IEC~60335
  & Límites alineados con la norma; consumo validado; trazabilidad metrológica
    pendiente de patrón \\
Cronograma & Tolerancia de 10 días
  & Cumplido: el hardware se demoró unos días por los envíos de componentes,
    dentro de la tolerancia, y el software se adelantó unas dos semanas
    (Anexo~\ref{anx:cronograma}) \\
Presupuesto & $\leq$ \SI{1030}{USD}
  & Cumplido: unos \SI{520}{USD} con contingencia
    (Tabla~\ref{tab:costos_detalle}) \\
\bottomrule
\end{tabularx}
\end{table}
